
\chapter{LITERATURE SURVEY AND OBJECTIVES} % Main chapter title
\label{ChapterLiteratureSurvey} % For referencing the chapter
% elsewhere, use \ref{Chapter1}

\section{Critical Review of Literature}
A growing prevalence of digital services has resulted in a need for
Terms of Service (ToS) agreements governing user-provider relations.
However, the complexity and length of such documents have given rise 
to a number of important research developments in the automated 
analysis of the contractual content itself. The paper extends the 
research on legal infomatics, Natural Language Processing (NLP)\index{NLP}, 
and knowledge graph-based reasoning\index{Knowledge Graph} into the 
domain of interdisciplnary issues. Among all the research, 
the CLAUDETTE framework is arguably the most important and has 
introduced automated identification of potentially unfair 
clauses\index{Clause Detection} included in ToS agreements, 
employing ML algorithms.
\\
The paper proved that it could offer better precision and recall 
in comparison to other classification techniques. 
Though it has come so far as a work of such significance to the 
field, the CLAUDETTE model was largely dependent on the 
classification of clauses and did not consider the contextual 
dependencies that could exist between clauses, i.e., 
it did not reason on complex structures of contracts.
Research in a similar area to contract law has revealed the consequences
on society in general in regards to unfair contract provisions and the
importance of intervention in such matters. As for unfair contract
term legislation, research in the legal area of unfair contracts revealed
the discrepancy between service providers and service consumers,
emphasizing the need for instruments that would promote service
provision and awareness among consumers.
\\
Although the above works give a good theoretical foundation they do not 
suggest any technological tools 
that can be automated for consumer use. Recent studies on the 
interface of AI and contract law suggest neural and machine learning 
technology for the automation of interpretation of contracts. 
The above studies suggest that AI can help in identifying unfair 
terms in contracts and can assist legal practitioners in interpreting 
contracts. However, it also suggest that AI is limited in terms of 
interpretation, content and context.
\\
Recent breakthroughs achieved using Retrieval-Augmented Generation 
(RAG)\index{RAG} and graph-based techniques have created new 
opportunities for dealing with these problems.
The use of Graph RAG has highlighted the importance of legal 
texts as structured texts that have hierarchical and relational relationships.
The application of Graph RAG has provided better contextual 
retrieval\index{Contextual Retrieval} compared to vector-based RAG methods,
though with a focus on legislative texts rather than private contract
agreements.
\\
Benchmarking studies on measuring large language models\index{LLM} for contract
review have also identified gaps in thinking over long context and clause
retrieval. Linguistic models may not identify information on large documents,
and hence the interest in hybrid models\index{Hybrid Retrieval}. Hybrid
architectures for applications such as financial analysis using Hybrid
Graph-RAG models\index{Hybrid Retrieval} that include vector similarity
search\index{Vector Similarity Search} and graph traversal\index{Graph Traversal} have also been proposed.
While promising, such techniques were not specifically developed for the analysis of legal documents, 
such as contracts, or for identifying clauses based on risks. 
To summarise, the literature reviewed above clearly shows that 
there has been substantial improvement in the automated techniques 
for contract analysis, as well as the outstanding issues, which include 
structural reasoning, explainability, and user-centered risk interpretability. 
This has greatly affected the proposed RiskWise framework.

\section{Literature Collection \& Segregation}
The studies evaluated in this study systematically identified
academic repositories like Google Scholar, arXiv, IEEE Xplore, and
reference tracing from relevant publications. 
The identified studies were grouped into thematic groups for analysis.
In the first category, studies on the detection of unfair clauses and 
legal texts' classification, such as the
CLAUDETTE framework, as well as other studies on machine learning, 
were identified.
\\
The studies provided a basic understanding of risk identification 
through the application of automation 12
techniques, though most of the studies were limited to classification.
There were also studies on legal and policy studies regarding unfair 
contract terms and consumer protection policies.
The studies provided domain knowledge on different problematic clauses.
The third category of studies included AI-based legal analytics and 
understanding of contracts, which discussed the development of neural language models.
The fourth category of studies included those related to 
Retrieval-Augmented Generation, Knowledge Graph,
and hybrid retrieval models that have come into the limelight in recent times.
\\
This study has been a methodological inspiration
for the integration of semantic and structural models in document
understanding. This thematic categorization allowed the
identification of complementary studies among research areas, which
also helped in identifying methodological gaps.

\section{Summary of Literature}
According to a literature review, automated legal document analysis
is progressively moving from traditional machine learning
classification techniques to more sophisticated language model-based
techniques. 
Though classification-based systems are geared towards recognizing 
particular clause patterns, there is a limit to their explanatory 
potential and a lack of consideration of clause interactions.
Though providing few computational results, the legal literature 
highlights the need to address clauses of unfair contract terms from a 
social perspective.
\\
Though AI-based systems have certain limitations related to context, 
they have a high potential for understanding text.
As a measure to improve context aggregation, there is a focus on Graph RAG, 
and its potential is emphasized, including structural representation and 
multi-hop reasoning.
However, these architectures have yet to undergo a thorough 
implementation in consumer-oriented ToS analysis.
Therefore, in conclusion, the literature indicates the feasibility of 
ToS analysis and the requirement for systems
comprising structural reasoning and semantic retrieval and risk 
assessment in an explainable manner.

\section{Identification of Gaps}
Based on the critical review and synthesis of the literature
surveyed, some research gaps were identified that form the basis for
the development of the proposed system.
\\
Firstly, it has been observed that
the existing systems for detecting unfair clauses, for the most part, rely on classification models that classify each clause individually.
Though it may have been successful in identifying the clauses of a particular kind, it may not always be possible to take into consideration the context of the document in which these clauses may be presented.
It may not always be possible to arrive at a reliable risk assessment depending on the relationship between clauses, supporting statements, and references, which may or may not be applicable.
\\
Secondly, it may not always be very transparent with regard to the results obtained by existing AI-based contract analysis tools.
Although these tools are successful in identifying clauses that may be potentially problematic, they are often not
very forthcoming about the rationale used for risk classification.
\\
Third, the conventional Retrieval-Augmented Generation models generally work directly with flat forms of texts.
This may lead to the problem of the "fragmentation of contextual information," where "semantically related clauses
that are spread out in different parts of the document may not be retrieved together." As such, "crucial cross-referenced
information may not be taken into account in the reasoning process."
\\
Fourthly, it has been shown that there is a limit in the way in which large scale models of language process complex relational dependencies in a document.
In agreements where clauses and relationships are interdependent, it has been shown that there is a difficulty in maintaining a coherent state of understanding in large-scale models of language.
\\
Finally, existing research on the domain of automated contract analysis tools has been found to be mostly focused on
the solutions that are meant for legal professionals and enterprise applications.
These applications have been found to be used by domain experts and may require output in technical or legal
formats that may not be easily understandable by a non-expert.
This shows that there is a problem related to a lack of user-friendly tools that can interpret complex contract
risks in an understandable format for a non-expert.



\section{Need for Research/Project}
The gaps identified above indicate that there is a noticeable need
for a smart system, providing semantic and structural representations
on ToS documents and explaining results to non-expert users. Such a
system would be beneficial for increasing transparency in digital
agreements, promoting informed consent, and increasing consumer
awareness. Additionally, the increasing usage trend of generative AI
can be leveraged to develop hybrid architectures targeting
limitations of current techniques. This project seeks to bridge the
gap between automated detection and human-understandable explanations
of contractual risks by integrating knowledge graph reasoning with
RAG-based language models.

\section{Problem Statement}
Despite the existence of Terms of Service contracts on digital
platforms, users are not prepared to understand complex contractual
terms and clauses that could negatively affect their rights, privacy,
financial commitments, and service access terms. 
\\
The contracts are complex, lengthy, and use legal terminology that is not understood by any layman,
which results in the acceptance of terms of the contract by the users without proper consideration.
In terms of the difference between acceptance of users and awareness of users, this is one of the biggest
risks in terms of decision-making in a digital ecosystem.
In terms of the analysis of contracts, though a few automatic methods have been proposed, there are
several such methods that have some serious flaws.
\\
Some are largely based on clause-level classification methods that are unable to reason in context and fail to
establish the relationship between distributed contractual clauses.
Some methods provide simple and generic models of summarization based on language models that do not
specifically point to potentially risky clauses or explain them in a manner that is interpretable to establish risks.
Therefore, the need for a system that can mediate with contextual awareness and structured reasoning as well as
user-centered interpretation, is still required.
\\
Hence the objective of the present paper is to develop an intelligent and transparent system that can analyse and comprehend the Terms of Service documents in an automated manner, as well as detect the critical clauses of the contract in a manner that can be easily interpreted by the user.
The solution to the problem would be achieved through the development of a hybrid retrieval and reasoning paradigm to assist the use of semantic retrieval, knowledge graph construction, and generative language models.


\section{Objectives}
The main aims and objectives of the project are as follows:

\begin{enumerate}
  \item Design and develop a hybrid retrieval system that combines
    vector similarity search and knowledge graph traversal for legal
    document analysis.
  \item Design and develop an automated system that can extract
    structured knowledge representations from unstructured Terms of
    Service documents.
  \item Design and develop a system that can identify and classify
    potentially risky contractual clauses using context-aware reasoning.
  \item Design and develop a system that can provide explainable
    natural language interpretations of identified risks.
  \item Design and develop an interactive web-based interface that
    can allow users to upload documents and perform automated risk analysis.
\end{enumerate}

\section{Limitations of the Study}

The study has a number of limitations. It is based on language
models, whose performance may depend on the capability of the model
and the complexity of the documents. The system is mainly suited for
English-language documents and is not applicable in multilingual
settings. In addition, the study has used qualitative methods of
evaluation because the study is exploratory in nature, and
large-scale quantitative benchmarking is not within the scope of the study.
Resource constraints related to local model execution and the
variability of LLM outputs are also limitations of the current system.
